\documentclass[12pt]{article}
\usepackage[utf8]{inputenc}
\usepackage[english,russian]{babel}
\usepackage{amssymb, amsmath, ragged2e, fancyhdr}
\linespread{1.0}

\pagestyle{empty}

\begin{document}

\setcounter{equation}{2}
\pagestyle{fancy}
\fancyhf{}
\lhead{УДК 519.715}

\renewcommand{\headrulewidth}{0pt}
\begin{center}
ПОСТРОЕНИЕ ГРАФА ПЕРЕХОДОВ\\
ПОСЛЕДОВАТЕЛЬНОСТНОЙ СХЕМЫ\\
С ПРИМЕНЕНИЕМ SAT-РЕШАТЕЛЯ\\
Иванов Иван Иванович\\
кандидат технических наук, доцент\\
доцент, Томский государственный университет\\
Россия, Томск\\
Петров Петр Петрович\\
студент, Томский государственный университет\\
Россия, Томск\\
\end{center}

\hspace{0.5cm}
Аннотация. Рассматривается подход к построению графа\\
переходов последовательностной схемы. Исследуется метод,\\
основанный на использовании SAT-решателя. Для построения графа\\
переходов определяются возможные переходы между состояниями\\
последовательной схемы и используются предварительные\\
вычисления, основанные на троичном и двоичном моделировании,\\
значительно сокращающие объем вычислений. Также\\
рассматривается построение на основе графа переходов\\
последовательностей, обеспечивающих заданные переходы схемы.\\
Компьютерные эксперименты показали эффективность\\
предложенного метода построения графа переходов с применением\\
SAT-решателя.

\hspace{0.5cm}
Ключевые слова: граф переходов, троичное моделирование,\\
SAT-решатель, последовательностная схема, переходная\\
последовательность.\\


\newline
\hspace{0.5cm}
В работе рассматривается подход к построению графа переходов\\
синхронной последовательностной схемы. Построение графа,\\
основанное на использовании SAT-решателя, исследуется подробно.\\
Представлены результаты компьютерных экспериментов для\\
предложенного метода построения графа переходов, применяющего\\
SAT-решатель. Также кратко рассматривается решение задачи\\
построения последовательности входных векторов, обеспечивающей\\
переход схемы в одно из состояний заданного множества, по графу\\
переходов.\\

\newpage
\hspace{0.5cm}
Рассмотрим синхронную последовательностную схему с \textit{n}\\
входами, \textit{m} выходами и \textit{p} элементами памяти (триггерами).\\
\textit{X} = $\{\textit{{x_1, ..., x_n}}\}$ - \textnormal{входные переменные схемы}, \textit{Y} = $\{\textit{{y_1, ..., y_m}}\}$ - \textnormal{ее}\\
\textnormal{выходные переменные}, \textit{Z} = $\{\textit{{z_1, ..., z_p}}\}$ - \textnormal{внутренние переменные\\
схемы.}

\hspace{0.5cm}
\textnormal{Назовем} графом переходов последовательностной схемы\\
\textnormal{ориентированный граф, у которого вершины сопоставлены\\
состояниям схемы и есть дуга из вершины} i \textnormal{в вершину} j \textnormal{тогда и\\
только тогда, когда в схеме существует одношаговый переход из\\
состояния, соответствующего вершине} i, \textnormal{в состояние,\\
соответствующее вершине} j, \textnormal{при каких-либо значениях входных\\
переменных}.

\hspace{0.5cm}
\textnormal{На рисунке 1 представлена комбинационная составляющая} С\\
\textnormal{последовательностной схемы. При построении графа переходов\\
рассматривается только та часть схемы, которая необходима для\\
получения функций переходов, выходы схемы не рассматриваются.\\
Поэтому структурное описание комбинационной составляющей,\\
используемой для получения графа переходов, упрощается (рисунок\\
2).}\\




\end{document}
